\documentclass[12pt, a4paper]{article}

% ── Packages ──────────────────────────────────────────────────────────────────
\usepackage[margin=2.5cm]{geometry}
\usepackage{times}
\usepackage[T1]{fontenc}
\usepackage[utf8]{inputenc}
\usepackage{amsmath, amssymb}
\usepackage{graphicx}
\usepackage{hyperref}
\usepackage{booktabs}
\usepackage{enumitem}
\usepackage{xcolor}
\usepackage{listings}
\usepackage{caption}
\usepackage{subcaption}
\usepackage{float}
\usepackage{titlesec}
\usepackage{abstract}
\usepackage{parskip}
\usepackage{microtype}
\usepackage{cite}
\usepackage{url}
\usepackage{multirow}
\usepackage{array}
\usepackage{tabularx}

% ── Hyperref Configuration ─────────────────────────────────────────────────────
\hypersetup{
    colorlinks=true,
    linkcolor=blue!70!black,
    citecolor=blue!70!black,
    urlcolor=blue!70!black,
    pdftitle={Design and Development of an Intelligent ArXiv Research Paper Ingestion Agent with RAG-Based Retrieval},
    pdfauthor={},
}

% ── Listings (code) ────────────────────────────────────────────────────────────
\lstset{
    basicstyle=\ttfamily\small,
    breaklines=true,
    frame=single,
    backgroundcolor=\color{gray!8},
    keywordstyle=\color{blue!70!black}\bfseries,
    commentstyle=\color{green!50!black},
    stringstyle=\color{red!70!black},
    numbers=left,
    numberstyle=\tiny\color{gray},
    numbersep=6pt,
    captionpos=b,
}

% ── Section Formatting ─────────────────────────────────────────────────────────
\titleformat{\section}{\large\bfseries}{\thesection.}{0.5em}{}
\titleformat{\subsection}{\normalsize\bfseries}{\thesubsection.}{0.5em}{}
\titleformat{\subsubsection}{\normalsize\itshape}{\thesubsubsection.}{0.5em}{}

% ── Title Block ────────────────────────────────────────────────────────────────
\title{%
    \textbf{Design and Development of an Intelligent ArXiv Research Paper
    Ingestion Agent with RAG-Based Retrieval} \\[0.4em]
    \large Monthly Progress Report --- March 2026
}

\author{
    Omar S.\\[0.2em]
    \small Supervised Research Project \\
    \small \textit{Department of Computer Science} \\[0.4em]
    \small \texttt{April 2026}
}

\date{}

% ══════════════════════════════════════════════════════════════════════════════
\begin{document}
\maketitle
\thispagestyle{empty}

% ── Abstract ──────────────────────────────────────────────────────────────────
\begin{abstract}
This report documents one month of iterative design and development of an
intelligent agent system for the ingestion, processing, and retrieval of
ArXiv research papers. Beginning from a conceptual prototype, the system
evolved through three weekly iterations into a containerised, RAG-enabled
pipeline comprising a natural-language query interface, a PDF extraction
layer, an embedding store, and a vector-similarity retrieval component.
The report covers the architectural decisions made at each iteration,
evaluates the PDF parsing tools examined (PDFPlumber, PyMuPDF, and
Docling), describes the chunking and embedding strategies adopted, and
identifies planned directions toward a metadata-driven advanced RAG
framework. The technology stack is grounded in LangChain, LangGraph,
Sentence Transformers, ChromaDB, and the GROQ inference API.

\noindent\textbf{Keywords:} ArXiv; LLM agents; Retrieval-Augmented Generation;
PDF extraction; vector databases; LangGraph; ChromaDB; document ingestion.
\end{abstract}

\vspace{1em}
\hrule
\vspace{1.5em}

% ── TOC ───────────────────────────────────────────────────────────────────────
\tableofcontents
\newpage

% ══════════════════════════════════════════════════════════════════════════════
\section{Introduction}
\label{sec:intro}

The exponential growth of scientific literature hosted on repositories such
as ArXiv necessitates intelligent tools capable of going beyond keyword
search: tools that understand natural-language intent, locate semantically
relevant documents, and synthesise coherent answers from their content.
Retrieval-Augmented Generation (RAG) \cite{lewis2020rag} provides a
principled framework for grounding large language model (LLM) outputs in
external, verifiable documents, thereby mitigating hallucination while
preserving the fluency of generative responses.

This project aims to design and implement such a system --- an
\emph{ArXiv Research Paper Ingestion Agent} --- that operates as an
end-to-end pipeline from user query to cited, contextualised response.
The development follows an iterative software-engineering methodology
with weekly enhancement cycles.

The remainder of this report is structured as follows. Section~\ref{sec:background}
reviews the relevant background and related work. Section~\ref{sec:arch}
presents the overall system architecture. Sections~\ref{sec:week1}--\ref{sec:week3}
detail the work accomplished in each of the three development weeks.
Section~\ref{sec:eval} evaluates the tools and strategies explored.
Section~\ref{sec:future} outlines future directions, and
Section~\ref{sec:conclusion} concludes.

% ══════════════════════════════════════════════════════════════════════════════
\section{Background and Related Work}
\label{sec:background}

\subsection{Retrieval-Augmented Generation}

Retrieval-Augmented Generation (RAG), introduced by Lewis et al.\
\cite{lewis2020rag}, combines a parametric language model with a
non-parametric dense retriever to condition generation on retrieved
passages. The architecture has since been extended in multiple directions,
including iterative retrieval \cite{shi2024iter}, graph-based
knowledge integration \cite{edge2024graphrag}, and
metadata-filtered retrieval. The present project implements a standard
(Simple RAG) pipeline as its baseline before advancing toward a
metadata-driven variant.

\subsection{Scientific Document Processing}

Processing scientific PDFs presents well-known challenges: multi-column
layouts, mathematical notation, figures, and tables require robust parsers
that go beyond na\"{i}ve text extraction. Prior work has evaluated tools
such as GROBID \cite{lopez2009grobid}, PDFMiner, and more recently
Docling \cite{docling2024} for structured extraction from scholarly
documents. Evaluation of these tools under practical conditions forms a
central contribution of this report.

\subsection{Agentic Frameworks for Information Retrieval}

The emergence of LLM-based agent frameworks --- LangChain \cite{langchain},
LangGraph \cite{langgraph} --- enables the construction of multi-step
reasoning pipelines where tools (e.g., an ArXiv API wrapper, a vector
store, a PDF parser) are invoked dynamically based on model decisions.
This paradigm underpins the agent design described herein.

% ══════════════════════════════════════════════════════════════════════════════
\section{System Architecture}
\label{sec:arch}

\subsection{High-Level Overview}

The system adopts a modular \emph{fan-out / fan-in} pattern orchestrated
by a LangGraph state machine. At a high level the pipeline consists of
five stages:

\begin{enumerate}[leftmargin=*, label=\textbf{\arabic*.}]
    \item \textbf{Query Parser} --- Translates a natural-language user
          query into structured search parameters (title, author, abstract
          keywords, category codes, date ranges) compatible with the
          ArXiv Atom API.
    \item \textbf{Source Router} --- Routes incoming requests to the
          appropriate retrieval source: the live ArXiv API or a locally
          ingested document corpus.
    \item \textbf{Document Ingestion Pipeline} --- Downloads and processes
          PDF documents via a three-stage sub-pipeline:
          \emph{extraction} $\rightarrow$ \emph{chunking}
          $\rightarrow$ \emph{embedding and storage}.
    \item \textbf{Retrieval Engine} --- Performs vector-similarity search
          against the embedded corpus stored in ChromaDB.
    \item \textbf{Response Aggregator} --- Synthesises a coherent,
          cited answer from the retrieved chunks using an LLM.
\end{enumerate}

\subsection{Technology Stack}

Table~\ref{tab:stack} summarises the principal technologies employed.

\begin{table}[H]
    \centering
    \caption{Technology stack employed in the ArXiv Ingestion Agent.}
    \label{tab:stack}
    \begin{tabularx}{\linewidth}{@{}l l X@{}}
        \toprule
        \textbf{Component} & \textbf{Technology} & \textbf{Role} \\
        \midrule
        LLM Inference & GROQ API & Cloud-hosted, low-latency LLM serving \\
        LLM (local) & Ollama & On-device model hosting for offline use \\
        Agent Framework & LangGraph + LangChain & Agent orchestration and tool binding \\
        PDF Extraction & PyMuPDF (\texttt{pymupdf4llm}) & Primary text extraction from PDFs \\
        Embedding Model & Sentence Transformers (HuggingFace) & Dense vector representation \\
        Vector Store & ChromaDB & Local vector database with similarity search \\
        ArXiv Interface & ArXiv Python API / feedparser & Metadata and full-text retrieval \\
        Containerisation & Docker & Reproducible deployment \\
        IDE / Dev & VS Code & Development environment \\
        \bottomrule
    \end{tabularx}
\end{table}

% ══════════════════════════════════════════════════════════════════════════════
\section{Week 1: Prototyping and Architectural Design}
\label{sec:week1}

\subsection{Objectives}

The first development week was dedicated to establishing a conceptual
baseline and producing a working prototype agent capable of answering
queries via the ArXiv metadata API. Discussions with the supervising
researcher provided key requirements that shaped the subsequent design.

\subsection{Architectural Workflow Design}

An initial architectural workflow was drafted based on the fan-out/fan-in
pattern described in Section~\ref{sec:arch}. The core data flow was
defined as:

\[
  \texttt{Query Parser} \;\longrightarrow\; \texttt{ArXiv Search}
  \;\longrightarrow\; \texttt{Agent(s)}
  \;\longrightarrow\; \texttt{Response Aggregator}
\]

This decomposition enforces a clean separation of concerns and allows
individual components to be upgraded independently --- a design principle
that proved valuable in subsequent iterations.

\subsection{ArXiv Repositories Agent}

The prototype agent translates free-form natural-language queries into
structured ArXiv API search parameters. Specifically, the query parser
identifies and maps the following fields:

\begin{itemize}[leftmargin=*]
    \item \textbf{Title} (\texttt{ti:}) --- keywords extracted from the
          query targeting the paper title.
    \item \textbf{Author} (\texttt{au:}) --- author names where mentioned.
    \item \textbf{Abstract} (\texttt{abs:}) --- semantic keywords for
          abstract matching.
    \item \textbf{Category} (\texttt{cat:}) --- ArXiv subject classification
          codes (e.g., \texttt{cs.AI}, \texttt{cs.CL}).
    \item \textbf{Date range} (\texttt{submittedDate:}) --- temporal
          constraints expressed in the query.
\end{itemize}

Responses returned to the user include: paper titles, author lists,
abstracts, publication dates, and direct hyperlinks to the ArXiv landing
pages.

\subsection{Initial Codebase}

The codebase was structured as a modular Python project. Core modules
implemented during this week included:

\begin{itemize}[leftmargin=*]
    \item \textbf{Query parsing module} --- LLM-driven extraction of
          structured fields from a natural-language query.
    \item \textbf{ArXiv API integration layer} --- wrapper around the
          ArXiv Atom feed using \texttt{feedparser} and \texttt{requests}.
    \item \textbf{LLM tool definitions} --- LangChain tool objects
          exposing search and metadata functions to the agent.
    \item \textbf{Response formatter} --- template-based rendering of
          retrieved metadata into human-readable output.
\end{itemize}

Initial testing was performed against both a locally hosted LLM via
Ollama and the GROQ cloud inference API, validating that the agent
architecture is backend-agnostic.

% ══════════════════════════════════════════════════════════════════════════════
\section{Week 2: Document Processing Pipeline}
\label{sec:week2}

\subsection{Objectives}

The second iteration extended the system beyond metadata retrieval to
full-document ingestion. The key deliverables were: a version-controlled
repository, a PDF extraction component, an embedding pipeline, and a
stored vector index ready for retrieval.

\subsection{Repository and Version Control}

A GitLab repository was created, and the existing codebase was pushed
with a structured \texttt{README} documenting the agent architecture,
component descriptions, and setup instructions. This established the
collaborative baseline for subsequent development.

\subsection{PDF Extraction Component}

PyMuPDF (\texttt{pymupdf4llm}) was selected as the primary PDF
extraction library following a comparative evaluation (detailed in
Section~\ref{sec:eval}). The extraction module:

\begin{itemize}[leftmargin=*]
    \item Ingests a PDF document from a local path or a URL.
    \item Extracts the textual content page-by-page.
    \item Exports the structured text to a markup file (\texttt{.md})
          for downstream processing, using an agentic tool to persist the
          intermediate artefact.
\end{itemize}

\subsection{Agent Routing}

A \emph{source router} node was added to the LangGraph state machine.
This node classifies each incoming request as one of:

\begin{description}[leftmargin=*]
    \item[\texttt{arxiv\_live}] Queries answered via live ArXiv API calls.
    \item[\texttt{local\_document}] Queries answered from the locally
          ingested and embedded document corpus.
\end{description}

The routing decision is driven by an LLM classifier that analyses the
user intent and the available data sources.

\subsection{Document Chunking}

After extraction, text is split into overlapping fixed-length chunks
using LangChain's \texttt{RecursiveCharacterTextSplitter}. While this
simple chunking strategy is functional, it does not respect semantic
boundaries. Semantic chunking --- which aligns chunk boundaries with
topic transitions --- is identified as a near-term improvement
(Section~\ref{sec:future}).

\subsection{Embedding and Vector Storage}

Sentence-level embeddings are generated using a Sentence Transformer
model from the HuggingFace \texttt{sentence-transformers} library.
The resulting dense vectors, together with the original text chunks
and document-level metadata, are persisted in a ChromaDB collection
for efficient approximate nearest-neighbour search.

% ══════════════════════════════════════════════════════════════════════════════
\section{Week 3: Retrieval Pipeline and Containerisation}
\label{sec:week3}

\subsection{Objectives}

The third iteration closed the loop from ingestion to answer generation
by implementing the retrieval pipeline and completing the Docker
containerisation for reproducible deployment.

\subsection{Retrieval Pipeline}

The retrieval component integrates similarity search directly into the
LangGraph agent workflow. Upon receiving a query, the pipeline:

\begin{enumerate}[leftmargin=*, label=\textbf{\arabic*.}]
    \item Embeds the query using the same Sentence Transformer model used
          at ingestion time, ensuring vector-space consistency.
    \item Issues a $k$-nearest-neighbour query against the ChromaDB
          collection.
    \item Ranks and filters results by cosine similarity score.
    \item Packages the top-$k$ chunks as context for the LLM response
          generator.
\end{enumerate}

The result is a functional Simple RAG architecture \cite{lewis2020rag}
providing grounded, document-backed answers with low hallucination risk.

\subsection{Containerisation}

The full agent stack --- including the Python application, ChromaDB, and
the Sentence Transformer inference server --- was containerised using
Docker. A multi-service \texttt{docker-compose} configuration was defined,
enabling a single-command deployment of the entire pipeline. This
significantly reduces environment-related friction and aligns with
best practices for reproducible research software.

\subsection{System Diagrams}

Detailed workflow diagrams, abstract architectural charts, and system
flow diagrams were produced this week to document the internal state
transitions of the LangGraph agent, the data flow through the ingestion
pipeline, and the interaction between the retrieval and generation
components.

\subsection{Literature Reviews Conducted}

Three areas of related literature were surveyed this week:

\begin{itemize}[leftmargin=*]
    \item \textbf{Metadata extraction methods}: MeExtract and related
          structured-extraction pipelines from scientific PDFs.
    \item \textbf{RAG pipeline comparison}: Simple RAG graphs versus
          graph-augmented retrieval \cite{edge2024graphrag}.
    \item \textbf{Advanced RAG}: Iterative and adaptive retrieval
          strategies relevant to the planned metadata-driven extension.
\end{itemize}

% ══════════════════════════════════════════════════════════════════════════════
\section{Evaluation of PDF Parsing Tools}
\label{sec:eval}

Three PDF parsing libraries were evaluated during Week~2 on a test corpus
of ArXiv papers spanning diverse layouts (single-column, double-column,
figures-heavy, and tables-heavy documents). Table~\ref{tab:parsers}
summarises the comparative findings.

\begin{table}[H]
    \centering
    \caption{Comparative evaluation of PDF parsing libraries.}
    \label{tab:parsers}
    \begin{tabularx}{\linewidth}{@{}l X X X@{}}
        \toprule
        \textbf{Library} & \textbf{Strengths} & \textbf{Weaknesses} & \textbf{Selected?} \\
        \midrule
        \textbf{PDFPlumber} &
            Simple API; accurate for single-column, text-only PDFs &
            Struggles with multi-column layouts, embedded images, and tables; low recall on complex documents &
            No \\[0.4em]
        \textbf{PyMuPDF} (\texttt{pymupdf4llm}) &
            Fast; good handling of multi-column layouts; LLM-optimised markdown output; lightweight &
            Limited table reconstruction fidelity in some edge cases &
            \textbf{Yes} \\[0.4em]
        \textbf{Docling} &
            Best overall structural extraction; handles tables and figures well &
            Heavy resource footprint; longer processing time; occasional disordered output on complex layouts &
            Candidate for future use \\
        \bottomrule
    \end{tabularx}
\end{table}

PyMuPDF was selected as the primary extractor due to its favourable
balance of speed, extraction quality, and its native LLM-oriented
markdown output format (\texttt{pymupdf4llm}). Docling remains a
candidate for a future high-fidelity extraction mode where processing
time is less critical.

% ══════════════════════════════════════════════════════════════════════════════
\section{Future Directions}
\label{sec:future}

The following enhancements are planned for the next development phase:

\subsection{Semantic Chunking}

The current fixed-length chunking strategy is known to split semantically
coherent passages at arbitrary boundaries, degrading retrieval precision.
Semantic chunking methods --- which detect topic transitions and align
chunk boundaries accordingly --- will be investigated and implemented.
A dedicated Jupyter notebook will be developed to evaluate and visualise
chunk quality under both strategies.

\subsection{Metadata-Driven RAG}

The system will be extended from a Simple RAG baseline to a
\emph{metadata-driven RAG} architecture. This involves:

\begin{itemize}[leftmargin=*]
    \item Storing rich document-level metadata (title, authors, venue,
          year, ArXiv category, citation count) alongside the chunk
          embeddings.
    \item Enabling metadata-filtered retrieval queries that constrain the
          search space prior to embedding-based similarity ranking.
    \item Integrating structured metadata extraction pipelines (e.g.,
          MeExtract-style approaches) to populate these fields
          automatically.
\end{itemize}

\subsection{Microsoft GraphRAG Integration}

Microsoft GraphRAG \cite{edge2024graphrag} constructs a knowledge graph
from the document corpus and enables community-level summarisation,
offering qualitatively different retrieval semantics from dense vector
search. An exploratory integration is planned to evaluate whether
graph-based context retrieval improves answer quality for multi-hop
reasoning queries.

\subsection{API-Based Storage and External Database}

Current deployment stores all data locally. Future iterations will
expose a REST API for document ingestion and retrieval, and evaluate
managed vector database services (e.g., Qdrant) for scalability and
additional filtering features.

\subsection{Parser Refinement}

A systematic Jupyter notebook-based evaluation will compare PDFPlumber,
PyMuPDF, and Docling across a larger benchmark corpus, providing
quantitative extraction quality metrics (character error rate, structural
recall) to inform future parser selection.

% ══════════════════════════════════════════════════════════════════════════════
\section{Planned Literature Survey}
\label{sec:lit}

The following papers have been identified for close review in the
upcoming development period:

\begin{enumerate}[leftmargin=*, label={[\arabic*]}]
    \item \textit{MIRIAD: Augmenting LLMs with Millions of Medical
          Query-Response Pairs} \cite{miriad2025} --- domain-specific
          QA dataset construction at scale; relevance to curating
          ArXiv-derived QA pairs.
    \item \textit{Agent-based Multimodal Information Extraction for
          Nanomaterials} \cite{agentnano2025} --- multi-modal agentic
          extraction; relevance to handling figures and equations in
          scientific PDFs.
    \item \textit{PDFChatAnnotator: A Human-LLM Collaborative
          Multi-Modal Data Collection Tool for PDF-Format Catalogs}
          \cite{pdfchat2024} --- human-in-the-loop annotation of
          PDF content; relevance to ground-truth construction for
          extraction evaluation.
\end{enumerate}

% ══════════════════════════════════════════════════════════════════════════════
\section{Conclusion}
\label{sec:conclusion}

Over the course of three iterative development weeks, the ArXiv Research
Paper Ingestion Agent evolved from a conceptual prototype into a
containerised, end-to-end RAG pipeline. Key contributions include:

\begin{itemize}[leftmargin=*]
    \item A modular, LangGraph-orchestrated agent architecture supporting
          both live ArXiv metadata queries and locally ingested PDF
          corpora.
    \item A comparative evaluation of three PDF extraction libraries,
          resulting in the adoption of PyMuPDF as the primary parser.
    \item An embedding and vector-storage pipeline using Sentence
          Transformers and ChromaDB.
    \item A functional similarity-search retrieval component completing
          the Simple RAG loop.
    \item A fully containerised deployment via Docker.
\end{itemize}

The system currently provides grounded, document-backed conversational
responses via the GROQ inference API. Planned extensions toward
semantic chunking, metadata-driven retrieval, and graph-augmented RAG
will progressively increase both the depth and precision of the
system's research assistance capabilities.

% ══════════════════════════════════════════════════════════════════════════════
\section*{Acknowledgements}

The author thanks the supervising researcher for guidance on system
requirements and for productive discussions that shaped the architectural
decisions documented in this report.

% ══════════════════════════════════════════════════════════════════════════════
\bibliographystyle{ieeetr}
\begin{thebibliography}{99}

\bibitem{lewis2020rag}
P.\ Lewis, E.\ Perez, A.\ Piktus, F.\ Petroni, V.\ Karpukhin,
N.\ Goyal, H.\ K\"{u}ttler, M.\ Lewis, W.\ Yih, T.\ Rockt\"{a}schel,
S.\ Riedel, and D.\ Kiela,
``Retrieval-Augmented Generation for Knowledge-Intensive NLP Tasks,''
in \textit{Advances in Neural Information Processing Systems (NeurIPS)},
vol.~33, pp.~9459--9474, 2020.

\bibitem{shi2024iter}
W.\ Shi, S.\ Min, M.\ Yasunaga, M.\ Seo, R.\ James, M.\ Lewis,
L.\ Zettlemoyer, and W.\ Yih,
``REPLUG: Retrieval-Augmented Language Model Pre-Training,''
in \textit{Proc.\ NAACL}, 2023.

\bibitem{edge2024graphrag}
D.\ Edge, H.\ Trinh, N.\ Cheng, J.\ Bradley, A.\ Chao,
A.\ Mody, S.\ Truitt, and J.\ Larson,
``From Local to Global: A Graph RAG Approach to Query-Focused
Summarization,''
\textit{arXiv preprint arXiv:2404.16130}, 2024.

\bibitem{lopez2009grobid}
P.\ Lopez,
``GROBID: Combining Automatic Bibliographic Data Recognition and
Term Extraction for Scholarship Publications,''
in \textit{Proc.\ 13th European Conference on Digital Libraries
(ECDL)}, pp.~473--474, 2009.

\bibitem{docling2024}
IBM Research,
``Docling: An Open-Source Document Conversion Toolkit,''
\textit{GitHub Repository}, \url{https://github.com/DS4SD/docling}, 2024.

\bibitem{langchain}
H.\ Chase,
``LangChain: Building Applications with Large Language Models,''
\textit{GitHub Repository}, \url{https://github.com/langchain-ai/langchain}, 2022.

\bibitem{langgraph}
LangChain AI,
``LangGraph: Stateful, Multi-Actor Applications with LLMs,''
\textit{GitHub Repository}, \url{https://github.com/langchain-ai/langgraph}, 2024.

\bibitem{miriad2025}
[Authors],
``MIRIAD: Augmenting LLMs with Millions of Medical Query-Response Pairs,''
\textit{arXiv preprint arXiv:2506.06091}, 2025.

\bibitem{agentnano2025}
[Authors],
``Agent-based Multimodal Information Extraction for Nanomaterials,''
\textit{npj Computational Materials}, 2025.
\url{https://www.nature.com/articles/s41524-025-01674-7}

\bibitem{pdfchat2024}
[Authors],
``PDFChatAnnotator: A Human-LLM Collaborative Multi-Modal Data
Collection Tool for PDF-Format Catalogs,''
in \textit{Proc.\ ACM IUI}, 2024.
\url{https://dl.acm.org/doi/10.1145/3640543.3645174}

\end{thebibliography}

\end{document}
